% !TEX root = ./paper.tex
\section{Discussion}\label{sec:discussion}
\subsection{Running Symbolic Execution Multiple Times with Small Budget}

\begin{figure}
    \centering
    \includegraphics[width=\textwidth]{figure/parasuit_vs_klee.pdf}
    \caption{Efficiency of $\ourtool$ in branch coverage compared to a single long-run of KLEE.}
    \label{fig:branchdiff}
\end{figure}

\textcolor{blue}{$\ourtool$ attempts various parameter values by running symbolic execution multiple times within the total time budget, using the small budget for each iteration.} We evaluated whether this approach is more efficient than running symbolic execution once with a long budget from the perspective of covered branches and cache efficiency. $\kleeone$, which ran for the same 24 hours as $\ourtool$, and $\kleedouble$, which doubles the budget, were defined as baselines. The same five benchmarks as in the experiments were used.
% ($\xorriso$, $\gcal$, $\grep$, $\gawk$, $\combine$)

Figure~\ref{fig:branchdiff} demonstrates the efficiency of running symbolic execution multiple times with a small budget. $\ourtool$ covered 83.8\% more branches than $\kleeone$, and even when running KLEE twice the budget, it only explored 59.8\% of the branches that $\ourtool$ covered. Furthermore, the set of branches covered by $\ourtool$ is close to the superset of the branch set by the two baselines. Within the same budget, $\ourtool$ missed only an average of 31.6 branches that $\kleeone$ covered, yet uniquely covered 1,858 branches--approximately 58 times more--and even outperformed $\kleedouble$ by covering 14.4 times more unique branches. Additionally, when comparing $\kleeone$ and $\kleedouble$, $\kleedouble$ missed an average of 174.2 branches that $\kleeone$ covered, which is 5.5 times more than the number of missed branches compared to $\ourtool$.

$\ourtool$ also showed acceptable performance in terms of cache efficiency. We calculated the percentage of how many queries can be determined to be satisfiable with the cache without using a solver. $\ourtool$ solved 78\% of the queries through the cache, showing similar cache efficiency to $\kleeone$ and $\kleedouble$, which achieved 81.3\% and 75.2\%, respectively. These experimental results demonstrate that trying various parameter sets for a short budget is more efficient than using a fixed parameter set for a long budget.


\subsection{Comparison with the Non-adaptive Approach}
\begin{figure}
    \centering
    \begin{subfigure}{0.47\textwidth}
        \includegraphics[width=\textwidth]{figure/param_diffrence.pdf}
        \caption{The difference in the parameter selection.}
    \end{subfigure}
    \begin{subfigure}{0.47\textwidth}
        \includegraphics[width=\textwidth]{figure/value_diffrence.pdf}
        \caption{The difference in the value sampling.}
    \end{subfigure}
    \caption{The differences in parameter tuning results between $\ourtool$ and $\symtuner$.}
    \label{fig:diffrences}
\end{figure}

\textcolor{blue}{$\ourtool$ tunes program-adaptively which parameters to use and what values to assign to those parameters.} We analyzed how this approach differs from previous techniques and what causes the performance differences. To do this, we directly compared $\ourtool$ with SymTuner~\cite{symtuner}, the state-of-the-art parameter tuning method for symbolic execution. We considered both the parameter selection and value sampling stages using the same five benchmarks as in the previous results. 

Figure~\ref{fig:diffrences}-(a) shows that $\ourtool$ tunes optimal parameters that previous techniques fail to tune. $\ourtool$ automatically selects an average of 49.6 parameters for each benchmark, while SymTuner manually selects 20 parameters that are equally tuned regardless of the benchmark. $\ourtool$ selects parameters that improve the performance of a specific benchmark through trial-and-error, but SymTuner misses an average of 72.6\% (36 parameters) that improve performance when tuned. This result causes the performance difference between $\ourtool$ and SymTuner in the parameter selection stage.

Figure~\ref{fig:diffrences}-(b) illustrates $\ourtool$ tries various values compared to previous techniques and converges to an optimal value. All three integer parameters analyzed in Figure~\ref{fig:sampled}-(b) are also selected by SymTuner. When comparing the values assigned to parameters by each technique, $\ourtool$ covers the entire range explored by SymTuner across all benchmarks, while SymTuner explores only an average of 12.4\% of the range that $\ourtool$ explores. Since SymTuner follows an algorithm that selects one of five predefined values manually as shown in Figure~\ref{fig:params}-(b), the actual range it explores is even less than 12.4\%. In addition, the best-performing value for most parameters is outside the range defined by SymTuner. For example, in the case of ${\sf sym}$-${\sf stdin}$, which was selected in all benchmarks, SymTuner selects one of the values 4, 8, 12, 16, and 20. However, the optimal values of ${\sf sym}$-${\sf stdin}$ discovered by $\ourtool$ are the values that SymTuner fails to assign in four benchmarks except $\gawk$(5): $\xorriso$(542), $\gcal$(107), $\grep$(29), and $\combine$(40). Based on these results, performance differences between $\ourtool$ and the previous technology also occur at the parameter value sampling stage.


%%% Local Variables: 
%%% mode: latex
%%% TeX-master: "paper"
%%% End:
